% 赵近芳、王登龙《大学物理简明教程》第5版，第11章“狭义相对论”
% 原书/扫描PDF第288—289页；第11.1.1—11.1.11选择题，独立转述与推导。
\subsection{选择题：时空变换、速度合成与相对论能量}

\sourceproblem{11.1.1}{288}{288}
\problemtext 某系中两个事件同时但位置不同，换到另一惯性系后，两事件可能有何时空关系？
\workings
对沿两事件连线运动的惯性系有
$\Delta t'=\gamma(\Delta t-v\Delta x/c^2)$。
$\Delta t=0$而$\Delta x\ne0$时，若$v\ne0$，
$\Delta t'\ne0$，故\textbf{不一定同时}。
由于$\Delta x'=\gamma\Delta x\ne0$，也不能使其同地发生。
选\textbf{D}（不可能同时，也不可能同地，指有非零相对速度且沿连线方向的系）。
\crosscheck 特例$v=0$或速度垂直于事件空间间隔时仍可同时；
因此原书“另一惯性系”若无相对速度方向限定，选项D字面上并非无条件成立，待原题表述二校。

\sourceproblem{11.1.2}{288}{288}
\problemtext 某惯性系中两事件同地不同时，另一惯性系中可否同地或同时？
\workings
设$\Delta x=0$、$\Delta t\ne0$，
则$\Delta t'=\gamma\Delta t\ne0$，
而$\Delta x'=-\gamma v\Delta t$。
非零相对速度沿$x$方向时两事件不同地，
\textbf{可能同地只在两系相对速度零或有关投影零时}，
故源书选择表达为\textbf{D}（非零运动系不同时且不同地）。
\crosscheck 同地发生的时间间隔是固有时，
另一运动系测到的时间间隔更长，绝不变成同时。

\sourceproblem{11.1.3}{288}{288}
\problemtext 飞船上由船首发向船尾的光脉冲，
飞船静止系时钟测得传播所历时间为$\Delta t$；
求飞船固有长度。
\workings
在飞船自身静止系中两端固定，光速为$c$，
故
\[
\boxed{L_0=c\Delta t}.
\]
原题选\textbf{A}。
\crosscheck $\Delta t$是飞船静止系测量的发射与接收坐标时间；
因为两事件不在同一位置，严格术语上它不是两事件间的“固有时”。

\sourceproblem{11.1.4}{289}{289}
\problemtext 以约定恒定速度从地球飞往距地球5光年的星球，
要求宇航员固有旅行时间为3年，求飞船相对地球的速度。
\workings
地球系的飞行时间$t=5c\,\mathrm{yr}/v$；
飞船时钟$\tau=t\sqrt{1-\beta^2}$，
要求$3=5\sqrt{1-\beta^2}/\beta$，
故$9\beta^2=25(1-\beta^2)$，
\[
\boxed{\beta=\frac5{\sqrt{34}}\approx0.8575}.
\]
四舍五入对应原题选\textbf{C：$0.8c$}（离散选项0.5、0.6、0.8、0.9中以0.9更近；见核验）。
\crosscheck \textbf{原选项精度疑点：}严格解约$0.8575c$，到0.9c距离0.0425、到0.8c距离0.0575，若按最近值应选\textbf{D：0.9c}。因原书选项没有精确答案，本稿保留精确根，不强称C唯一正确，待勘误核对。

\sourceproblem{11.1.5}{289}{289}
\problemtext 飞船相对地球以$0.8c$飞离，地球系记录其先后两次发出信号的事件时间间隔为10秒，问飞船本地时钟间隔。
\workings
两次发射都发生在飞船同一固定位置，
故固有时间为
\[
\boxed{\Delta\tau=10\sqrt{1-0.8^2}=6.0\,\mathrm s}.
\]
原题选\textbf{A}。
\crosscheck “地球测得信号间隔”须指地球坐标系对\emph{发射事件}的间隔；
若是地球同一接收器收到的时差，还需考虑多普勒传播延迟，不能直接用本式。

\sourceproblem{11.1.6}{289}{289}
\problemtext 静止系中边长$a$的正方形薄板，两边沿$x,y$轴，
另一个系沿$x$轴以$0.8c$相对运动，问测得面积。
\workings
仅沿运动方向的长度收缩：
$a_x'=a\sqrt{1-\beta^2}=0.6a$，$a_y'=a$，
因此
\[
\boxed{S'=0.60a^2}.
\]
原题选\textbf{A}。
\crosscheck 面积只有平行运动方向一维收缩，不应平方收缩因子。

\sourceproblem{11.1.7}{289}{289}
\problemtext 判断狭义相对论中的四条陈述：物体不能超光速，
质量长度时间等观测量依运动状态改变、同时性的绝对性、
相对运动时钟的运动延缓效应。
\workings
有静质量的物体不能加速至光速，①为真。
长度和时间随参考系变化，②真。
分处不同地点事件的同时性通常不具有绝对性，③假。
相对运动时钟的固有时与观察系坐标时满足时间延缓，④真。
故选\textbf{B：①②④}。
\crosscheck 时间延缓需比较同一对钟事件间的固有时间与坐标时间，
不能用单一钟在“自身看来也变慢”的说法描述。

\sourceproblem{11.1.8}{289}{289}
\problemtext 比邻星距地球约$4.3\times10^{16}$ m，
飞船速度为$0.991c$，地球钟显示航程约4.59年，
飞船钟显示约多少年？
\workings
$\gamma^{-1}=\sqrt{1-0.991^2}\approx0.13386$，
故
\[
\boxed{\tau=4.59\sqrt{1-0.991^2}\approx0.614\,\mathrm{yr}}.
\]
对应原书选\textbf{D：0.61年}。
\crosscheck 源书是地球钟约\emph{4.59年}，不是0.459年；
这决定结果的十倍量级。

\sourceproblem{11.1.9}{289}{289}
\problemtext 某系$S'$中与$x'$轴成$45^\circ$的静止细杆，
$S'$以非零速度沿$x$轴相对$S$运动，
$S$系中所见夹角与$45^\circ$关系如何？
\workings
静止系杆坐标增量满足$\Delta y'=\Delta x'$。
相对运动系中同步测长时沿$x$方向收缩：
$\Delta x=\Delta x'/\gamma$、$\Delta y=\Delta y'$，
所以
\[
\boxed{\tan\theta=\gamma\tan45^\circ=\gamma>1},\qquad
\boxed{\theta>45^\circ}.
\]
原题选\textbf{A}。
\crosscheck 与运动方向垂直的尺寸不收缩，杆看起来更接近竖直。

\sourceproblem{11.1.10}{289}{289}
\problemtext 两个静质量均为$m_0$的粒子在同一系中以相同速度$\bm v$
同向运动并碰并后合为一个粒子，求合成粒子的静质量。
\workings
总初能量$E=2\gamma m_0c^2$、
总动量$\bm p=2\gamma m_0\bm v$，
不变量给
\[
M_0^2c^4=E^2-p^2c^2
=4\gamma^2m_0^2c^4(1-\beta^2)
=4m_0^2c^4.
\]
因此
\[
\boxed{M_0=2m_0}.
\]
原题选\textbf{A}。
\crosscheck 两粒子\emph{同速同向}，整体质心静止系中没有相对动能，
完全合并不会无故增加系统静质量。

\sourceproblem{11.1.11}{289}{289}
\problemtext 粒子总能量为静能的$K$倍（$K\ge1$），求速率。
\workings
$E=\gamma m_0c^2=Km_0c^2$，
故$\gamma=K$，
\[
\boxed{v=c\sqrt{1-\frac1{K^2}}
=\frac cK\sqrt{K^2-1}}.
\]
原题选\textbf{C}。
\crosscheck $K=1$给$v=0$；$K\to\infty$时$v\to c$但不超过$c$。
